\chapter*{Περίληψη}
\addcontentsline{toc}{chapter}{Περίληψη}
\label{ch:abstract-el}

Η επιλογή \tl{top-k}, δηλαδή ο εντοπισμός των \tl{k} μεγαλύτερων στοιχείων ενός
συνόλου, εμφανίζεται σε πλήθος εφαρμογών και επαναλαμβάνεται αρκετά συχνά ώστε το
κόστος της να είναι αισθητό. Η πρόσφατη ενσωμάτωση μονάδων νευρωνικής
επεξεργασίας στην ίδια ψηφίδα με τον επεξεργαστή καθιστά πρακτικό το ερώτημα αν
συμφέρει η ανάθεση τέτοιων υπολογισμών σε επιταχυντή.

Η εργασία υλοποιεί δύο αλγορίθμους με θεμελιωδώς διαφορετική συμπεριφορά, την
περικομμένη διτονική ταξινόμηση που είναι ανεξάρτητη από τα δεδομένα και την
επιλογή με τοπικούς σωρούς που εξαρτάται από αυτά, σε τρεις αρχιτεκτονικές:
επεξεργαστή με διανυσματικές επεκτάσεις \tl{AVX-512}, μονάδα γραφικών με
\tl{CUDA} και μονάδα νευρωνικής επεξεργασίας \tl{AMD XDNA}. Και
οι έξι υλοποιήσεις μετρήθηκαν στον ίδιο παραμετρικό χώρο, σε δύο ρητά διακριτά
πεδία μέτρησης, με καταγραφή χρόνου, ενέργειας και όγκου διακινούμενων δεδομένων
και με τοποθέτηση κάθε εκτέλεσης σε οροφή μετρημένη στο ίδιο μηχάνημα.

Τα αποτελέσματα δείχνουν ότι το πρόβλημα περιορίζεται από την κίνηση δεδομένων σε
κάθε αρχιτεκτονική, ότι η κατάταξη των δύο αλγορίθμων ανατρέπεται με τον λόγο
$n/k$ και όχι με την αρχιτεκτονική και ότι η επιλογή του ορίου μέτρησης
ανατρέπει πλήρως την κατάταξη των αρχιτεκτονικών. Ενεργειακά, η μονάδα
νευρωνικής επεξεργασίας εμφανίζει τη χαμηλότερη ισχύ και ταυτόχρονα την υψηλότερη
κατανάλωση, καθώς η διάρκεια κυριαρχεί του γινομένου. Καμία εκτέλεση σε
επιταχυντή δεν είναι βέλτιστη κατά \tl{Pareto}: για μικρό \tl{k} η ανάθεση είναι
μετρήσιμα χειρότερη από τον υπολογισμό στον επεξεργαστή.

\vspace{1em}

\textbf{Λέξεις-κλειδιά:} επιλογή \tl{top-k}, δίκτυα διτονικής ταξινόμησης,
ετερογενής υπολογισμός, μονάδα νευρωνικής επεξεργασίας, \tl{MLIR-AIE},
διανυσματοποίηση \tl{SIMD}, ενεργειακή αποτίμηση, μοντέλο \tl{roofline}

\selectlanguage{english}

\chapter*{Abstract}
\addcontentsline{toc}{chapter}{Abstract}
\label{ch:abstract-en}

Top-k selection, the task of identifying the k largest elements of a set, appears
across a wide range of applications and is repeated often enough that its cost is
noticeable. The recent integration of neural processing units into the same
die as the processor makes the question of whether offloading such
computations to an accelerator is worthwhile a practical one.

This work implements two algorithms with fundamentally different behaviour,
truncated bitonic sorting which is data-oblivious and map-reduce selection with
local heaps which is data-dependent, on three architectures: a processor with
AVX-512 vector extensions, a GPU programmed in CUDA, and an AMD XDNA neural
processing unit. All six
implementations were measured over the same parameter space, under two explicitly
distinct measurement scopes, recording time, energy and the volume of data moved
and placing every run against a roofline measured on the same machine.

The results show that the problem is limited by data movement on every
architecture, that the ranking of the two algorithms is overturned by the ratio $n/k$
rather than by the architecture and that the choice of measurement boundary
completely overturns the ranking of the architectures. In energy terms the neural
processing unit exhibits the lowest power draw and simultaneously the highest
consumption, since duration dominates the product. No accelerator run is
Pareto-optimal: for small k, offloading is measurably worse than computing on the
processor.

\vspace{1em}

\textbf{Keywords:} top-k selection, bitonic sorting networks, heterogeneous
computing, neural processing unit, MLIR-AIE, SIMD vectorisation, energy
efficiency, roofline model

\selectlanguage{greek}
